\begin{table}[tbp]
  \centering
  \caption{Desempeño en el conjunto de prueba, pobreza monetaria, 2013$\to$2016.}
  \label{tab:desempeno_pr}
  \footnotesize
  \setlength{\tabcolsep}{4pt}
  \begin{tabular}{l ccc ccc cc}
    \toprule
    & & & & \multicolumn{3}{c}{\textbf{Umbral F$_2$}} & \multicolumn{2}{c}{\textbf{Decil superior}} \\
    \cmidrule(lr){5-7}\cmidrule(lr){8-9}
    \textbf{Modelo} & \textbf{AUC-PR} & \textbf{IC95\%} & \textbf{AUC-ROC} & \textbf{Rec.} & \textbf{Prec.} & \textbf{F$_1$} & \textbf{Prec.} & \textbf{Rec.} \\
    \midrule
    \multicolumn{9}{l}{\textit{Sin ingreso ni gasto (Modelo B)}} \\
    \quad \emph{Elastic net} (principal) & 0.435 & 0.40--0.48 & 0.739 & 0.89 & 0.27 & 0.41 & 0.50 & 0.25 \\
    \quad Random Forest & 0.416 & 0.38--0.46 & 0.742 & 0.90 & 0.27 & 0.42 & 0.49 & 0.25 \\
    \quad XGBoost & 0.420 & 0.38--0.46 & 0.742 & 0.94 & 0.25 & 0.39 & 0.49 & 0.25 \\
    \quad LightGBM & 0.389 & 0.35--0.43 & 0.731 & 0.91 & 0.26 & 0.40 & 0.49 & 0.25 \\
    \quad HistGradientBoosting & 0.413 & 0.38--0.46 & 0.735 & 0.88 & 0.27 & 0.42 & 0.49 & 0.25 \\
    \addlinespace
    \multicolumn{9}{l}{\textit{Con ingreso y gasto (Modelo A)}} \\
    \quad \emph{Elastic net} & 0.433 & 0.39--0.47 & 0.747 & 0.88 & 0.28 & 0.43 & 0.51 & 0.26 \\
    \quad Random Forest & 0.422 & 0.39--0.47 & 0.765 & 0.93 & 0.28 & 0.43 & 0.48 & 0.24 \\
    \quad XGBoost & 0.431 & 0.39--0.48 & 0.767 & 0.95 & 0.26 & 0.41 & 0.50 & 0.25 \\
    \quad LightGBM & 0.418 & 0.38--0.46 & 0.758 & 0.90 & 0.28 & 0.43 & 0.52 & 0.26 \\
    \quad HistGradientBoosting & 0.431 & 0.39--0.47 & 0.763 & 0.91 & 0.28 & 0.43 & 0.51 & 0.26 \\
    \bottomrule
  \end{tabular}
  \begin{minipage}{0.95\textwidth}
    \vspace{4pt}
    \footnotesize \textit{Nota:} AUC-PR: área bajo la curva de
    \emph{precision}-\emph{recall}, estimada como \emph{average precision};
    un modelo que asigna el riesgo al azar obtiene 0.199, la
    proporción de hogares vulnerables, y 0.5 de AUC-ROC. IC95\%:
    \emph{bootstrap} de 2{,}000 remuestras por comunidad. Umbral F$_2$:
    métricas al umbral elegido por validación cruzada maximizando F$_2$.
    Decil superior: si un programa atiende al 10\% de los hogares con
    mayor riesgo predicho. Rec.: \emph{recall}; Prec.: \emph{precision}.
    3{,}191 hogares no pobres en 2013; una estimación por modelo.
    Fuente: cálculos propios.
  \end{minipage}
\end{table}